\begin{lemma}\label{lem:decic-relative-formal-F4}
Let $R$ be the integers of the unramified quadratic extension of
$\mathbb Q_2$, with residue field $k=\mathbb F_4$.
Let $X/R$ be a smooth proper curve, let $A=\operatorname{Jac}(X)$,
and let $h$ be an involution acting freely on a neighbourhood of
the cusps. Put $\alpha=h_*-1$ and $\pi:X\to Z=X/\langle h\rangle$.
Suppose that $Z/R$ is smooth near the images of these cusps.
Write a cuspidal divisor $\overline C$ of degree $d$ as
\[
              \overline C=\overline C_0+\pi^*\overline E,
              \qquad \deg\overline E=k_0,
\]
where the multiplicities of $\overline E$ are maximal on each
pair interchanged by $h$. Let $C$ be its cuspidal lift.
Suppose that $\alpha[D-C]$ is torsion for an effective divisor
$D\in X^{(d)}(R)$ with $\overline D=\overline C$.
If restriction of
\[
 U=(\mathcal C^2-1)\alpha^*H^0(X_k,\Omega^1)
\]
to $\overline C$ has rank $d-k_0$, then $D$ contains the fixed
divisor $C_0$.
\end{lemma}
\begin{proof}
We first record the formal-group calculation needed at $2$.
For a torsion point $P$ in the first formal filtration of $A(R)$,
write its coordinates as $t(P)=2u$. The second filtration is
torsion-free, by the formal logarithm, and $[2]P=0$.
If $H$ is the matrix of the quadratic terms of $[2]$, reduction
of $[2](2u)=0$ modulo $8$ gives
\[
                   \overline u=H\overline u^{[2]}.
\]
Squaring and using $\overline u^{[4]}=\overline u$ gives
$\overline u=HH^{[2]}\overline u$.
The linearized Cartier matrix is $H^{[2]}$, so $HH^{[2]}$
is the matrix of $\mathcal C^2$ on cotangents, with the dual
coordinate convention. Thus every form in
$(\mathcal C^2-1)H^0(A_k,\Omega^1)$ annihilates the first
formal tangent of $P$.

Apply this to $P=\alpha[D-C]$.
The first tangent of $D-C$ is therefore annihilated by the
restriction of $U$. Its kernel has dimension $k_0$.
The formal family
\[
                  T=C_0+\pi^*Z^{(k_0)}
\]
through $C$ has dimension $k_0$ and lies in the fibre of
$E\mapsto\alpha[E-C]$.
It is a closed formal immersion into $X^{(d)}$.
Indeed, $\pi$ is \'etale at the cusps. On each pair of local
branches, the two moving monic polynomials are the same, and
maximality of $\overline E$ leaves no fixed factor on at least
one branch. The moving coefficients can therefore be recovered
from that branch. This also proves tangent injectivity.
Its tangent space is consequently the kernel of the restriction
of $U$.

Since $T/R$ is formally smooth, we may choose $C'\in T(R)$
with $C'\equiv D\pmod4$. We have
$\alpha[D-C']=\alpha[D-C]$, because $\alpha\pi^*=0$.
This torsion point lies in the second formal filtration, and
hence is zero.
Finally, the rank of $\alpha^*$ restricted to $\overline C$
is at least $d-k_0$. Choose $d-k_0$ independent cotangent
coordinates of the map $E\mapsto\alpha[E-C]$.
Their common formal zero locus is smooth of dimension $k_0$.
It contains $T$, and the tangent map of $T$ to this locus is an
isomorphism. The formal implicit function theorem therefore
identifies this zero locus with $T$. Since $D$ lies in the full
fibre, $D\in T(R)$, and $C_0\leq D$.
\end{proof}
